\documentclass[conference,onecolumn]{IEEEtran}

\usepackage[T1]{fontenc}
\usepackage[utf8]{inputenc}
\usepackage[spanish,es-tabla]{babel}
\usepackage{lmodern}
\usepackage{microtype}
\usepackage{csquotes}
\usepackage{amsmath,amssymb}
\usepackage{booktabs}
\usepackage{float}
\usepackage{tabularx}
\usepackage{array}
\usepackage{longtable}
\usepackage{url}
\usepackage[hidelinks]{hyperref}
\usepackage[backend=biber,style=ieee,sorting=none]{biblatex}
\addbibresource{referencias.bib}
\AtBeginDocument{\spanishplainpercent}

\newcolumntype{Y}{>{\raggedright\arraybackslash}X}
\newcolumntype{L}[1]{>{\raggedright\arraybackslash}p{#1}}
\renewcommand{\IEEEkeywordsname}{Palabras clave}

\title{FORM-PC-01: Comparación de modelos de proceso de requisitos en escenarios reales}
\author{\IEEEauthorblockN{Arteaga Alava Danela Dayana (líder), Chavarría Cuenca Tahiny Mel,\\
Naranjo Alfonzo Anthony Alexander, Nivela Saavedra José Ricardo}
\IEEEauthorblockA{Grupo ACNN · Carrera de Software\\
Facultad de Ciencias de la Computación\\
Universidad Técnica Estatal de Quevedo, Ecuador\\
Ingeniería de Requerimientos -- ISR-401}}

\begin{document}
\maketitle

\begin{abstract}
% Responsable: Danela. Se redacta al final, cuando todas las secciones estén completas.
\end{abstract}

\begin{IEEEkeywords}
Ingeniería de requisitos, modelos de proceso, cascada, iterativo, ágil, framework de Pohl.
\end{IEEEkeywords}

\section*{Datos generales}

\begin{table}[H]
\centering\footnotesize
\begin{tabularx}{\textwidth}{@{}p{0.22\textwidth}Y@{}}
\toprule
\textbf{Campo} & \textbf{Detalle}\\
\midrule
Asignatura & Ingeniería de Requerimientos (ISR-401) · 4.\textsuperscript{o} nivel · Paralelo B\\
Grupo & ACNN\\
Integrantes & Arteaga Alava Danela Dayana (líder), Chavarría Cuenca Tahiny Mel, Naranjo Alfonzo Anthony Alexander, Nivela Saavedra José Ricardo\\
Docente & Gleiston Cicerón Guerrero Ulloa, PhD\\
Fecha de cierre & Sábado 7 de noviembre de 2026, 23:59\\
Repositorio & \url{https://github.com/darteagaa-boop/ISR401-GrupoACNN}\\
\bottomrule
\end{tabularx}
\end{table}

\section{Criterios de comparación}
% Responsable: Danela · Ronda 1 (sábado, hasta 12:00)
% Definir al menos 5 criterios de comparación, por ejemplo: estabilidad de los
% requisitos, criticidad del sistema, disponibilidad del cliente, tamaño del
% proyecto, regulación y documentación exigida. Escribir una definición por criterio.


\section{Caracterización de los modelos de proceso}
% Responsable: Tahiny · Ronda 1 (sábado, hasta 12:00)

\subsection{Modelo en cascada}
\label{sec:modelo-cascada}

El modelo en cascada organiza el trabajo en fases sucesivas, con productos
revisados y aprobados antes de avanzar a la siguiente etapa. En ingeniería
de requisitos, su característica principal es establecer una especificación
suficientemente completa y validada antes de usarla como base del diseño y
la implementación. Esto facilita planificar entregables y conservar la
trazabilidad, pero también puede hacer costoso descubrir requisitos omitidos
en etapas tardías \cite{swebok2024,pohl2025}.

\subsection{Fases principales}
\begin{enumerate}
  \item \textbf{Delimitación y viabilidad.} Identificar el problema, los
  objetivos, las partes interesadas, el alcance, las restricciones y los
  riesgos iniciales. \emph{Resultado:} alcance y plan de trabajo acordados.

  \item \textbf{Obtención y análisis de requisitos.} Recopilar las
  necesidades de los actores, revisar el funcionamiento del entorno,
  resolver contradicciones y distinguir requisitos funcionales, de calidad
  y restricciones. \emph{Resultado:} requisitos analizados y priorizados.

  \item \textbf{Especificación y validación.} Documentar los requisitos
  con identificadores y criterios de aceptación; comprobar su coherencia,
  factibilidad, verificabilidad y correspondencia con las necesidades del
  cliente. Una vez aprobados, forman una \emph{línea base} sometida a
  control de cambios. \emph{Resultado:} especificación validada
  \cite{pohl2025,ireb2026}.

  \item \textbf{Diseño de la solución.} Definir arquitectura, módulos,
  interfaces y decisiones técnicas a partir de la especificación aprobada.
  \emph{Resultado:} diseños revisados y relacionados con sus requisitos.

  \item \textbf{Implementación e integración.} Construir los componentes e
  integrarlos respetando las interfaces acordadas. \emph{Resultado:}
  versión integrada del sistema.

  \item \textbf{Verificación, validación y aceptación.} Ejecutar pruebas
  de componentes, integración, sistema y aceptación; contrastar los
  resultados con requisitos y criterios previamente definidos.
  \emph{Resultado:} evidencias de conformidad y registro de defectos.

  \item \textbf{Despliegue y mantenimiento.} Instalar, operar, corregir
  defectos y tramitar nuevas necesidades mediante solicitudes de cambio
  evaluadas por su impacto. \emph{Resultado:} sistema operativo con
  documentación y versiones controladas.
\end{enumerate}

Las fases describen una secuencia predominante, no la prohibición absoluta
de retroceder: una no conformidad puede exigir revisar una fase anterior,
pero el cambio debe evaluarse y documentarse.

\subsection{Supuestos del modelo}
Para que esta organización sea razonable, se parte de los siguientes
supuestos, que deben comprobarse en el proyecto concreto:
\begin{enumerate}
  \item Las necesidades fundamentales pueden conocerse y aclararse antes
  de comprometer el diseño detallado.
  \item Los requisitos centrales y las interfaces tendrán una estabilidad
  relativa durante la construcción.
  \item Los responsables pueden revisar y aprobar entregables en hitos
  definidos, aunque no participen todos los días.
  \item Es viable especificar criterios de aceptación y conservar
  evidencias de cumplimiento para cada requisito relevante.
  \item Los cambios posteriores se gestionarán mediante un procedimiento
  explícito de análisis de impacto, aprobación y actualización documental.
\end{enumerate}

\subsection{Condiciones de aplicación y límites}
\textbf{Es favorable} cuando el alcance está delimitado, existen interfaces
técnicas conocidas, hay exigencias de documentación y trazabilidad, las
revisiones formales son importantes y el cliente puede aprobar una línea
base antes del desarrollo \cite{swebok2024}.

\textbf{Es poco conveniente como proceso estrictamente secuencial} cuando
los objetivos cambian con frecuencia, se desconoce la viabilidad técnica,
se necesita experimentación continua o las personas usuarias solo pueden
aclarar sus expectativas al probar prototipos. En estos casos, conviene
considerar un enfoque iterativo o introducir una exploración técnica
previa a la aprobación de la línea base.



% Responsable: Naranjo Alfonzo Anthony Alexander
% FORM-PC-01 · Ronda 1: caracterización del modelo iterativo.
\subsection{Modelo iterativo}
\label{sec:modelo-iterativo}

El modelo iterativo plantea que los requisitos no tienen que quedar completamente
resueltos en un único intento. Se comienza con una comprensión inicial del
problema y se realizan ciclos de obtención, documentación, negociación y
validación para corregir ambigüedades y ampliar la especificación. Cada ciclo
aprovecha lo aprendido en el anterior; por eso es útil cuando existen
necesidades conocidas, pero todavía quedan detalles por aclarar
\cite{sommerville2016,pohl2025}.

A diferencia de una secuencia estrictamente lineal, una observación del cliente
puede llevar a revisar requisitos ya documentados. El producto de cada ciclo
puede ser una versión mejorada de la Especificación de Requisitos de Software
(ERS), un prototipo evaluado o un conjunto de requisitos acordados. Iterar
sobre requisitos no obliga, por sí solo, a entregar software funcionando en
cada ciclo; esa decisión depende de cómo se organice el desarrollo.

\subsubsection{Fases principales}

\begin{enumerate}
  \item \textbf{Comprensión inicial y delimitación.} Identificar el problema,
  los objetivos, los actores, las restricciones y los requisitos conocidos.
  También se seleccionan los aspectos con mayor incertidumbre.
  \emph{Resultado:} alcance preliminar y primera versión de requisitos.

  \item \textbf{Planificación de la iteración.} Elegir qué funcionalidades,
  dudas o riesgos se analizarán en el ciclo y acordar quiénes revisarán los
  resultados. \emph{Resultado:} objetivos y criterios de revisión del ciclo.

  \item \textbf{Obtención y documentación.} Realizar entrevistas,
  observaciones o revisión de prototipos para precisar necesidades; actualizar
  requisitos funcionales, de calidad y restricciones, identificando sus
  fuentes. \emph{Resultado:} ERS o modelos revisados.

  \item \textbf{Negociación y validación.} Contrastar los requisitos con
  las partes interesadas, resolver contradicciones, comprobar que sean claros
  y verificables, y registrar los acuerdos. \emph{Resultado:} requisitos
  aceptados y observaciones pendientes \cite{ireb2026,pohl2025}.

  \item \textbf{Evaluación y siguiente ciclo.} Revisar las observaciones,
  controlar las versiones, actualizar las prioridades y decidir si se
  necesita otra iteración o si ya existe una especificación suficiente para
  continuar. \emph{Resultado:} nueva línea base o cierre documentado.
\end{enumerate}

La validación y la gestión de cambios acompañan a todo el proceso, no se
reservan únicamente para el último paso. Esto mantiene la trazabilidad entre
las necesidades del cliente, las versiones de los requisitos y las decisiones
tomadas \cite{pohl2025}.

\subsubsection{Supuestos del modelo}

\begin{itemize}
  \item Existen partes interesadas disponibles para revisar los resultados
  en momentos acordados y dar respuestas a las dudas importantes.
  \item Los requisitos pueden precisarse progresivamente sin que el costo de
  cada corrección haga inviable el proyecto.
  \item Es posible dividir el problema en conjuntos de requisitos o asuntos
  prioritarios que puedan analizarse y validarse por ciclos.
  \item El equipo dispone de mecanismos para registrar versiones, decisiones,
  dependencias y solicitudes de cambio, evitando perder acuerdos anteriores.
  \item Se puede reservar tiempo para la retroalimentación y el retrabajo
  que surjan al descubrir necesidades no previstas.
\end{itemize}

\subsubsection{Condiciones de aplicación y límites}

\textbf{Es conveniente} cuando el alcance general se conoce, pero aún hay
incertidumbre sobre pantallas, procesos, reglas de negocio o integración con
otros sistemas. También ayuda cuando los usuarios pueden revisar prototipos o
versiones intermedias y es necesario conservar una especificación formal que
mejore en cada ciclo \cite{sommerville2016}.

\textbf{Es menos conveniente} si no existe ningún representante capaz de
validar avances, si las decisiones dependen de autorizaciones demasiado
lentas o si el proyecto no dispone de tiempo para revisar el trabajo. En
sistemas con interfaces críticas o restricciones contractuales, las
iteraciones deben mantener controles de trazabilidad, aceptación y análisis
de impacto; repetir ciclos no elimina esas obligaciones.

\textbf{Diferencia frente al enfoque ágil.} Ambos aceptan revisiones, pero
el iterativo puede centrarse en refinar la ERS y validar prototipos mediante
ciclos planificados. En el enfoque ágil, la elaboración de requisitos suele
integrarse de forma continua con la construcción y las pruebas de incrementos
pequeños del producto \cite{schon2017}.

% Responsable: José · Ronda 1 (sábado, hasta 12:00)
% Describir las fases, los supuestos y las condiciones de aplicación del modelo.
\subsection{Modelo ágil}

El enfoque ágil organiza el desarrollo de software de forma incremental y empírica, con participación continua de las partes interesadas para obtener retroalimentación y ajustar el producto. Según Schön y sus colaboradores, las actividades de ingeniería de requisitos (obtención, documentación, validación, negociación y gestión) no están claramente separadas: se repiten durante las iteraciones y la información se detalla conforme se necesita para el trabajo siguiente \cite{schon2017}.

\subsubsection{Fases y ciclo de trabajo}

En los métodos ágiles las secuencias de fases no son únicas y en este caso se usará la propuesta de Scrum escrita por Schwaber \cite{schwaber1997}:

\begin{enumerate}
	\item \textbf{Preparación (Pregame):} Incluye planificación de la entrega de los elementos que se conocieron en el backlog con las estimaciones de tiempos y costos.
	
	\item \textbf{Desarrollo (Game).} Considerando los requisitos, el tiempo, la calidad, el costo y las condiciones del entorno se desarrollan ciclos iterativos que permiten evolucionar al sistema.
	
	\item \textbf{Cierre (Postgame).} Se prepara el producto para su entrega mediante tareas como integración, pruebas del sistema y documentación final.
\end{enumerate}

En la fase de desarrollo, Schwaber describe cuatro actividades recurrentes dentro del sprint \cite{schwaber1997}:

\begin{itemize}
	\item \textbf{Desarrollar (Develop):} analizar, diseñar, implementar, probar y documentar los cambios necesarios para atender los requisitos del backlog.
	
	\item \textbf{Preparar la versión ejecutable (Wrap):} reunir los cambios en una versión ejecutable que permita mostrar lo implementado.
	
	\item \textbf{Revisar (Review):} presentar el trabajo, evaluar el progreso, resolver problemas y revisar riesgos. Pueden identificarse nuevos elementos del backlog.
	
	\item \textbf{Ajustar (Adjust):} incorporar al producto la información y los cambios derivados de la revisión.
\end{itemize}

Esta descripción corresponde a Scrum y se usó como base debido a que Scrum es un marco de trabajo ágil, pero no es una secuencia universal en las metodologías ágiles.

\subsubsection{Supuestos}

Turk, France y Rumpe identifican supuestos que permiten analizar la aplicabilidad de los procesos ágiles. Para esta comparación se consideran los siguientes \cite{turk2014}:

\begin{itemize}
	\item \textbf{Visibilidad del progreso.} La entrega de software funcional funciona como una ventana para que el cliente analice los avances de manera más limpia sin la necesidad de leer los informes concurrentes, pero esta característica se ve limitada cuando el sistema no muestra su evolución mediante una interfaz gráfica como lo puede ser un sistema embebido.
	
	\item \textbf{División del trabajo en iteraciones cortas.} Se supone que el producto puede descomponerse en incrementos pequeños, desarrollables y demostrables en periodos breves; las dependencias fuertes entre subsistemas pueden dificultar esta organización.
	
	\item \textbf{Disponibilidad del cliente.} La interacción con el cliente es clave en desarrollo ágil y es necesaria una buena comunicación y disposición entre el equipo y el cliente, pero las responsabilidades habituales del cliente pueden limitar esa disponibilidad y, por consiguiente, retrasar la retroalimentación.
	
	\item \textbf{Comunicación directa.} Los autores identifican el supuesto de que la conversación cara a cara es la forma más productiva de comunicación entre clientes y desarrolladores, y entre los propios desarrolladores.
\end{itemize}

\subsubsection{Condiciones de aplicación}

El enfoque ágil puede funcionar cuando el cliente tiene tiempo para colaborar, las dudas se resuelven oportunamente y el trabajo puede dividirse en incrementos que se puedan evaluar. Esa participación permite aclarar los requisitos y cambiar las prioridades según lo que se aprende durante el desarrollo \cite{schon2017,turk2014}.

Si el cliente tiene poca disponibilidad o hay dependencias técnicas que impiden dividir el trabajo, será necesario adaptar la forma de trabajar. También puede ser más difícil coordinarse cuando el equipo está en distintos lugares. En sistemas sin interfaz humana, el avance puede mostrarse mediante simulaciones, pruebas de aceptación o revisiones, como señalan Turk y sus colaboradores \cite{turk2014}. Estas dificultades deben analizarse en cada escenario antes de decidir si conviene usar un enfoque ágil.

\section{Análisis de los escenarios}
% Responsable: Danela
\subsection{Sistema de un GAD con contrato público}

\subsubsection{Caracterización}
El docente presentó el escenario sin detalle, por lo que se describen los supuestos usados para el análisis.

\begin{table}[H]
\centering\footnotesize
\caption{Caracterización del escenario del GAD}\label{tab:gad}
\begin{tabularx}{\textwidth}{@{}p{0.26\textwidth}Y@{}}
\toprule
\textbf{Aspecto} & \textbf{Descripción del escenario}\\
\midrule
Sistema & Se asume que es una plataforma web para que los ciudadanos realicen trámites municipales en línea, ya sean solicitudes, pagos de tasas y seguimiento de trámites, para que los funcionarios los atiendan.\\
Cliente & El GAD municipal, representado por un administrador del contrato designado por la institución.\\
Actores del proceso & Administrador del contrato y fiscalizador por parte del GAD, funcionarios de los departamentos que atienden trámites, ciudadanos usuarios, autoridades del GAD que aprueban el alcance, el equipo proveedor y los organismos de control del uso de fondos públicos.\\
Estabilidad de los requisitos (C1) & Alta, ya que se asume que los requisitos se fijan en los términos de referencia del proceso de contratación pública antes de firmar el contrato, y cualquier cambio requiere una modificación formal.\\
Criticidad (C2) & Media, porque una falla no pone en riesgo vidas, pero afecta la atención a los ciudadanos y el manejo de datos personales y pagos.\\
Disponibilidad del cliente (C3) & Baja, ya que los funcionarios tienen otras responsabilidades y las decisiones pasan por aprobaciones internas, por lo que las reuniones son programadas y poco frecuentes.\\
Tamaño y complejidad (C4) & Grande, porque involucra varios departamentos, muchos tipos de trámite y la integración con sistemas existentes del GAD.\\
Regulación y contrato (C5) & Alta, el proyecto se rige por un contrato público con alcance, plazo y presupuesto definidos, y debe cumplir la normativa de contratación pública y de protección de datos personales.\\
Documentación exigida (C6) & Alta, porque se necesitan especificaciones, actas de entrega y recepción e informes de fiscalización para pagar cada entregable y responder ante los organismos de control.\\
\bottomrule
\end{tabularx}
\end{table}

\subsubsection{Modelo elegido y justificación}
\textbf{Modelo elegido:} Cascada

Se elige el modelo en cascada por las siguientes razones del escenario:

\begin{enumerate}
\item Los requisitos quedan fijados en los términos de referencia antes de firmar el contrato, y cualquier cambio exige una modificación formal. La cascada parte justamente de requisitos estables definidos al inicio.
\item El contrato público define alcance, plazo y presupuesto cerrados, y los pagos se hacen contra entregables aprobados por el fiscalizador. Las fases secuenciales de la cascada, con un entregable revisado al final de cada una, se ajustan a esa forma de aceptar el producto.
\item La disponibilidad del cliente es baja y la cascada concentra la participación de los funcionarios en la obtención y validación inicial de los requisitos, y no depende de reuniones frecuentes en cada iteración.
\item Se exige mucha documentación formal para pagar cada entregable y responder ante los organismos de control, la cascada produce una especificación completa y documentos de cada fase que sirven como evidencia.
\end{enumerate}

\textbf{Modelos descartados:}

\textbf{Ágil:} Supone requisitos cambiantes y un cliente disponible en cada iteración para priorizar y validar, condiciones que no se cumplen en el GAD. Además, su documentación mínima no alcanza para las actas y los informes que exige el contrato.

\textbf{Iterativo:} Se adapta mejor al tamaño del proyecto, porque permite entregar por partes la integración de varios departamentos. Sin embargo, revisar los requisitos en cada iteración choca con un alcance fijado por contrato y con un cliente poco disponible, por lo que obtiene menor total que la cascada en la matriz.

\subsubsection{Riesgos del modelo elegido y mitigación}
% Ronda 3 (sábado, hasta 22:30): riesgos concretos y cómo mitigarlos.

% Responsable: Tahiny
\subsection{Sistema embebido de una planta}

\textbf{Alcance del caso.} Se asume que una planta industrial dedicada a la
fabricación y envasado de productos necesita
implementar un sistema embebido que permita
automatizar y supervisar el funcionamiento de su
línea de producción. Actualmente, algunos procesos
requieren intervención manual, lo que puede
ocasionar retrasos, errores operativos, desperdicio
de materiales y dificultades para detectar fallos
en la maquinaria.
El sistema estará basado en un controlador embebido industrial que
recibirá información de sensores de temperatura,
presión, proximidad y nivel, con el propósito de
monitorear las condiciones de operación en tiempo
real. A partir de estos datos, el controlador
accionará automáticamente equipos como motores,
bandas transportadoras, bombas y válvulas,
siguiendo una secuencia de funcionamiento
previamente establecida. Asimismo, contará con una interfaz de supervisión
que permitirá a los operarios visualizar el estado
de los equipos, consultar parámetros de
funcionamiento y recibir alertas ante situaciones
anormales. Se contempla el registro de eventos
y la ejecución de acciones de protección, como
detener de manera controlada determinados equipos
cuando se detecten fallos.

\textbf{Límites del alcance:}
El proyecto se limita al monitoreo, control y
automatización de la línea de producción, sin
incluir procesos administrativos, comerciales
o financieros de la planta.
Se asume que los requisitos principales de
funcionamiento y seguridad serán definidos antes
del desarrollo, debido a la criticidad de las
operaciones industriales.

\subsubsection{Caracterización}
\subsection{Cliente, actores y características del proyecto}
\begin{description}
  \item[Cliente y objetivo.] Se asume que el cliente es la administración
  de la planta y que busca supervisar la operación de una línea productiva,
  reducir fallos de operación y registrar estados de los equipos.

  \item[Actores.] Se asume la participación de una jefatura de producción
  (acepta objetivos operativos), operadores (usan las señales y alarmas),
  personal de mantenimiento (diagnostica fallos), especialistas en
  automatización (integran sensores y actuadores) y responsables de
  seguridad industrial (revisan peligros y medidas de protección).

  \item[Estabilidad de requisitos.] \textbf{Alta en el núcleo, bajo
  supuesto:} se presume que la función de control y las interfaces físicas
  pueden definirse antes de construir. Los ajustes de parámetros o la
  presentación de alarmas podrían cambiar y requieren gestión de cambios.

  \item[Criticidad.] \textbf{Alta, bajo supuesto:} una orden incorrecta a
  un actuador o una lectura defectuosa podría afectar la continuidad del
  proceso, provocar daños en equipos o generar peligros. Esta valoración
  exige revisar y probar los comportamientos ante fallos; no presupone
  que todo el software cumpla una función de seguridad certificada.

  \item[Contrato y regulación.] Se asume un contrato o acuerdo técnico con
  alcance, interfaces y criterios de aceptación definidos. La normativa
  obligatoria \textbf{no está indicada en el enunciado} y deberá
  determinarse según país, maquinaria y funciones implementadas. Si hay
  funciones eléctricas/electrónicas/programables relacionadas con la
  seguridad, corresponde evaluar la pertinencia de IEC 61508-1:2010,
  sin afirmar que sea automáticamente obligatoria
  \cite{iec61508_2010}.

  \item[Disponibilidad del cliente.] \textbf{Limitada, bajo supuesto:}
  se presume que operadores y técnicos trabajan por turnos. Serían
  posibles reuniones programadas para aclarar requisitos, revisar
  documentos y participar en pruebas de aceptación, pero no necesariamente
  reuniones diarias.

  \item[Tamaño e integración.] \textbf{Medio, bajo supuesto:} el alcance
  incluye firmware o lógica de control, sensores, actuadores, una interfaz
  de operación y pruebas de integración con el equipo físico. No se
  dispone de estimaciones reales de esfuerzo ni de componentes.
\end{description}

\subsubsection{Modelo elegido y justificación}
\textbf{Elección propuesta: modelo en cascada con verificaciones técnicas
tempranas y control de cambios.} La elección solo se mantiene si la
revisión inicial confirma los supuestos anteriores. Se justifica por:
\begin{enumerate}
  \item \textbf{Estabilidad del núcleo de control:} permitiría acordar
  previamente qué señales se leen, qué acciones se autorizan y cuáles son
  los límites de operación, disminuyendo cambios tardíos en interfaces.
  \item \textbf{Criticidad y trazabilidad:} permite establecer requisitos
  y criterios de prueba identificables, relacionarlos con el diseño y
  conservar actas y resultados para la aceptación.
  \item \textbf{Participación del cliente por hitos:} se adapta mejor a
  revisiones programadas con producción, mantenimiento y seguridad
  industrial que a la exigencia de decisiones operativas diarias.
\end{enumerate}

\textbf{Por qué no el iterativo como modelo principal:} sería preferible
si existieran dudas relevantes sobre sensores, interfaces o respuestas del
controlador que solo pudieran resolverse mediante prototipos. Bajo los
supuestos actuales, esa incertidumbre se considera manejable con ensayos
preliminares, sin redefinir todo el proceso por ciclos.

\textbf{Por qué no el ágil como modelo principal:} las posibles ventanas
reducidas de acceso a la planta y la necesidad de acordar previamente las
interfaces y evidencias de aceptación hacen menos conveniente depender de
retroalimentación continua. Esto no significa que las prácticas ágiles
sean incompatibles con sistemas embebidos.

\subsubsection{Riesgos del modelo elegido y mitigación}
% Ronda 3 (sábado, hasta 22:30): riesgos concretos y cómo mitigarlos.

% Responsable: Naranjo Alfonzo Anthony Alexander
% FORM-PC-01 · Rondas 1, 2 y 3: escenario, elección y riesgos.
\subsection{Aplicación móvil de un emprendimiento}
\label{sec:escenario-app}

\subsubsection{Caracterización}

El enunciado únicamente indica una aplicación móvil de un emprendimiento,
sin precisar qué vende, su presupuesto o las condiciones de desarrollo.
Por ello, se plantea un \textbf{escenario hipotético}: un pequeño
emprendimiento de productos artesanales quiere ofrecer su catálogo y recibir
pedidos desde una aplicación móvil. Los datos siguientes son supuestos de
análisis, no resultados de entrevistas con un negocio real.

\paragraph{Descripción y alcance}
Se asume que la aplicación permitirá a los clientes consultar productos,
ver precios y disponibilidad, preparar una solicitud de pedido y conocer su
estado. El propietario podrá actualizar el catálogo y confirmar o rechazar
pedidos. En la primera versión, el pago y la entrega se coordinarán fuera de
la aplicación. No se incluirán pasarela de pagos, facturación electrónica,
servicio propio de reparto ni integración con sistemas empresariales.

\paragraph{Cliente y actores del proceso}
El cliente es la persona propietaria del emprendimiento, quien establece las
prioridades y acepta las funcionalidades. Los compradores son usuarios
principales y ayudan a evaluar la facilidad de uso. Se asume que una persona
encargada de atender pedidos participa en la definición del flujo de
confirmación; en un negocio pequeño puede ser la misma propietaria. El
equipo de desarrollo obtiene, documenta, negocia y valida los requisitos.

\paragraph{C1. Estabilidad de los requisitos: baja}
Aunque el objetivo de mostrar productos y recibir pedidos está claro, se
supone que la organización del catálogo, el proceso de compra, las
notificaciones y la presentación de las pantallas cambiarán después de que
el propietario y los clientes prueben las primeras versiones.

\paragraph{C2. Criticidad del sistema: baja}
Una falla podría provocar pedidos perdidos, molestias a los clientes o una
venta no concretada, pero el escenario no incluye control de equipos
peligrosos ni procesamiento de pagos dentro de la aplicación. La valoración
baja no elimina la necesidad de proteger los datos de contacto y verificar
que cada pedido quede registrado correctamente.

\paragraph{C3. Disponibilidad del cliente: alta}
Se asume que la persona propietaria puede responder consultas por mensajería
durante la semana y revisar una demostración breve al cerrar cada ciclo. Se
programarán sesiones puntuales con compradores para recibir opiniones sobre
el catálogo y la solicitud de pedidos.

\paragraph{C4. Tamaño y complejidad: pequeño}
El alcance inicial se limita al catálogo, los pedidos, su estado y la
gestión básica por parte del propietario. Participan pocos tipos de usuario
y no se contemplan integraciones externas complejas. Se supone un equipo de
desarrollo reducido.

\paragraph{C5. Regulación y contrato: baja}
Se asume que el proyecto no depende de una contratación pública ni de un
contrato que congele todas las funciones desde el inicio. Las prioridades
podrán acordarse directamente con la persona propietaria. Esto no significa
que no existan obligaciones legales: antes de operar habrá que revisar las
que correspondan al tratamiento de datos personales y a la actividad
comercial del emprendimiento.

\paragraph{C6. Documentación exigida: baja}
Se considera suficiente mantener historias de usuario, criterios de
aceptación, un registro de decisiones, evidencias básicas de prueba y una
guía de uso. No se supone una auditoría contractual que requiera informes
extensos para aceptar cada función.

\begin{table}[H]
\centering\footnotesize
\caption{Caracterización del escenario hipotético de la aplicación móvil}
\label{tab:caracterizacion-app}
\begin{tabularx}{\textwidth}{@{}Yp{0.20\textwidth}@{}}
\toprule
\textbf{Criterio} & \textbf{Valoración}\\
\midrule
C1. Estabilidad de los requisitos & Baja\\
C2. Criticidad del sistema & Baja\\
C3. Disponibilidad del cliente & Alta\\
C4. Tamaño y complejidad & Pequeño\\
C5. Regulación y contrato & Baja\\
C6. Documentación exigida & Baja\\
\bottomrule
\end{tabularx}
\end{table}

\subsubsection{Modelo elegido y justificación}

\textbf{Modelo elegido: ágil, con entregas incrementales y revisiones
frecuentes.} La matriz de la Tabla~\ref{tab:matriz-app} le asigna 18 puntos,
frente a 14 del iterativo y 9 de cascada. La selección depende de que se
cumpla la disponibilidad del propietario planteada como supuesto.

\begin{enumerate}
  \item \textbf{Requisitos cambiantes:} las pantallas y el recorrido para
  hacer pedidos pueden ajustarse cuando los compradores usen prototipos o
  versiones iniciales. El enfoque ágil permite revisar y priorizar esas
  necesidades continuamente \cite{schon2017}.

  \item \textbf{Participación del propietario:} las revisiones breves y las
  respuestas frecuentes permiten decidir rápidamente qué mejorar o qué
  función tiene más valor para el negocio.

  \item \textbf{Alcance pequeño y divisible:} se puede comenzar mostrando
  un catálogo funcional, añadir la solicitud de pedidos y luego incorporar
  la consulta de estados. Cada entrega permite observar problemas reales
  sin esperar a terminar toda la aplicación.

  \item \textbf{Pocas restricciones contractuales y documentales:} el
  emprendimiento puede reajustar prioridades sin tramitar modificaciones
  formales extensas, manteniendo de todas maneras criterios de aceptación,
  pruebas y registro de cambios.
\end{enumerate}

\textbf{Por qué no cascada como modelo principal.} Aunque facilitaría
establecer de antemano una especificación y un plan, exigiría decidir desde
el inicio detalles de experiencia de usuario que todavía no se conocen. Si
las pruebas con clientes muestran que el flujo de pedidos es confuso,
corregirlo tardíamente podría causar retrabajo.

\textbf{Por qué no el iterativo como modelo principal.} También es una
alternativa válida, ya que permite refinar requisitos y evaluar prototipos.
Sin embargo, bajo los supuestos de alta disponibilidad del propietario y
funciones pequeñas que pueden ponerse a prueba, resulta más conveniente
integrar las revisiones de requisitos con la construcción y entrega
frecuente de funcionalidades. El iterativo sería una opción más fuerte si
las revisiones solo pudieran efectuarse en hitos separados o se exigiera
una ERS formal actualizada en cada ciclo \cite{sommerville2016,schon2017}.

\subsubsection{Riesgos del modelo elegido y mitigación}

Los riesgos siguientes corresponden al escenario supuesto; no son problemas
comprobados en un emprendimiento real.

\begin{enumerate}
  \item \textbf{Crecimiento continuo del alcance.} El propietario podría
  solicitar pagos en línea, reparto propio u otras funciones fuera del
  objetivo inicial. \emph{Mitigación:} mantener una lista priorizada de
  historias de usuario, separar el alcance inicial de las mejoras futuras
  y revisar el costo y el tiempo antes de aceptar cambios relevantes.

  \item \textbf{Requisitos confusos o decisiones contradictorias.}
  Al recibir comentarios frecuentes, el equipo podría modificar pantallas
  sin criterios claros. \emph{Mitigación:} definir criterios de aceptación
  para cada historia, registrar decisiones y validar el flujo de pedido con
  el propietario antes de cerrar una entrega.

  \item \textbf{Pérdida de disponibilidad del cliente.} Aunque se supone
  alta participación, las actividades diarias del negocio pueden impedir
  revisar una versión. \emph{Mitigación:} acordar de antemano reuniones
  breves, nombrar a quien pueda responder dudas y no dar por aprobada una
  función que depende de una decisión pendiente.

  \item \textbf{Errores introducidos al agregar funciones.} Una nueva
  pantalla o actualización podría afectar la creación y el seguimiento de
  pedidos. \emph{Mitigación:} realizar pruebas del flujo completo de pedido
  en cada entrega, comprobar el registro de estados y probar la aplicación
  en dispositivos de distintas características.

  \item \textbf{Protección de datos y documentación descuidadas.} La
  rapidez de las entregas podría llevar a almacenar más información de la
  necesaria o a omitir registros de cambios. \emph{Mitigación:} recoger
  solo datos indispensables, establecer permisos de acceso, verificar
  requisitos de privacidad y mantener historias, decisiones y resultados
  de prueba vinculados a cada versión.
\end{enumerate}

% Responsable: José
\subsection{Plataforma de una cooperativa}

\subsubsection{Caracterización}

El enunciado únicamente menciona una plataforma de una cooperativa, sin especificar sus funciones ni condiciones de desarrollo. Para realizar la comparación, se plantea un caso hipotético de una cooperativa de ahorro y crédito. Las características y restricciones descritas son supuestos de trabajo, no datos obtenidos de una organización real.

\paragraph{Descripción de la plataforma}

Se supone que la cooperativa necesita una plataforma web para que sus socios consulten saldos y movimientos, realicen transferencias entre cuentas de la misma institución y presenten solicitudes de crédito. Los empleados podrán revisar las solicitudes y registrar su seguimiento; la aprobación del crédito continuará a cargo del personal autorizado.

La plataforma se conectará con un sistema existente que mantiene las cuentas y registra las operaciones. No sustituirá ese sistema ni calculará directamente los saldos. Quedan fuera del alcance inicial las transferencias a otras instituciones, la aplicación móvil y la aprobación automática de créditos.

\paragraph{Cliente y actores del proceso}

El cliente será la cooperativa, representada por su gerencia. Se consideran los siguientes actores:

\begin{itemize}
	\item \textbf{Gerencia:} acuerda el alcance, las prioridades y la aceptación del producto.
	
	\item \textbf{Responsables de operaciones y crédito:} explican los procedimientos, las reglas del negocio y las excepciones que debe contemplar la plataforma.
	
	\item \textbf{Socios:} aportan necesidades de uso y participan en la revisión de consultas, transferencias y solicitudes.
	
	\item \textbf{Responsable de tecnología:} proporciona información sobre el sistema existente y sus interfaces.
	
	\item \textbf{Responsable de seguridad y cumplimiento:} revisa las necesidades de protección de datos, permisos y registro de operaciones.
	
	\item \textbf{Equipo de desarrollo:} obtiene, documenta y analiza los requisitos, identifica conflictos y prepara los incrementos para su validación.
\end{itemize}

\paragraph{C1. Estabilidad de los requisitos: media}

Se supone que las operaciones principales y las reglas para registrar movimientos ya están definidas en el sistema existente. Sin embargo, los detalles de las pantallas, el seguimiento de solicitudes y las notificaciones necesitan aclararse con empleados y socios. Por ello, existe una base estable, pero parte de los requisitos puede ajustarse durante las revisiones.

\paragraph{C2. Criticidad del sistema: alta}

La plataforma permitirá acceder a información financiera e iniciar transferencias. Una falla podría exponer datos, duplicar una operación o mostrar información incorrecta. Por esta razón, el escenario requiere comprobar los permisos de acceso, la integridad de las operaciones y la correspondencia de lo mostrado con los registros del sistema existente.

\paragraph{C3. Disponibilidad del cliente: media}

Se supone que la cooperativa designará a un representante de operaciones para una reunión semanal y para responder dudas entre reuniones. Los responsables de crédito, seguridad y tecnología participarán en revisiones programadas. Los socios estarán disponibles para sesiones puntuales de validación, pero no para colaborar diariamente.

Esta disponibilidad permite recibir retroalimentación periódica, aunque algunas decisiones tendrán que esperar la intervención del responsable correspondiente.

\paragraph{C4. Tamaño y complejidad: mediano}

Se plantea un equipo de cinco personas para desarrollar una plataforma con tres grupos principales de funciones: consultas, transferencias internas y solicitudes de crédito. Participan usuarios con distintos permisos y existe una integración con el sistema de cuentas.

El alcance se considera mediano por estar limitado a esas funciones. No obstante, la integración y el manejo de operaciones financieras añaden complejidad técnica y requieren pruebas específicas.

\paragraph{C5. Regulación y contrato: alta}

Como supuesto del caso, la cooperativa exigirá revisiones de seguridad y cumplimiento antes de autorizar el uso de funciones que manejen datos u operaciones financieras. Las obligaciones normativas concretas deberán identificarse y verificarse antes de establecer la línea base de requisitos.

También se supone un contrato con presupuesto y fecha de entrega acordados, organizado por entregables. Los cambios que afecten alcance, costo o plazo deberán registrarse y aprobarse. Estas restricciones exigen controlar las modificaciones, aunque permiten revisar los requisitos dentro de los límites acordados.

\paragraph{C6. Documentación exigida: alta}

Se supone que la cooperativa solicitará una especificación de requisitos, reglas de negocio, criterios de aceptación, descripción de interfaces y registros de cambios. También exigirá trazabilidad entre los requisitos y sus pruebas, junto con evidencias de validación y aprobación de los entregables.

Esta documentación permitirá comprobar qué se acordó, por qué se modificó y con qué resultados se aceptó cada función.

\begin{table}[H]
	\centering
	\footnotesize
	\caption{Caracterización del escenario hipotético de la cooperativa}
	\label{tab:escenario-cooperativa}
	\begin{tabularx}{\textwidth}{@{}Yp{0.22\textwidth}@{}}
		\toprule
		\textbf{Criterio} & \textbf{Valoración} \\
		\midrule
		C1. Estabilidad de los requisitos & Media \\
		C2. Criticidad del sistema & Alta \\
		C3. Disponibilidad del cliente & Media \\
		C4. Tamaño y complejidad & Mediano \\
		C5. Regulación y contrato & Alta \\
		C6. Documentación exigida & Alta \\
		\bottomrule
	\end{tabularx}
\end{table}

Estas valoraciones describen el escenario hipotético. Se utilizarán como base para comparar los modelos en cascada, iterativo y ágil, sin asumir de antemano que alguno es el más adecuado.

\subsubsection{Modelo elegido y justificación}

\textbf{Modelo elegido: iterativo con desarrollo incremental y revisiones documentadas.}

Se propone construir versiones sucesivas, revisar los requisitos con la cooperativa y ampliar las funciones de forma controlada. Se distingue entre iterar, que permite revisar y mejorar lo trabajado, e incrementar, que añade funcionalidad. Ambas prácticas se combinan en la propuesta para este caso \cite{sommerville2016}.

La elección se apoya en las siguientes razones del escenario:
\begin{enumerate}
    \item \textbf{Hay requisitos que todavía necesitan aclararse.} Las operaciones principales tienen una base definida, pero las pantallas, las notificaciones y el seguimiento de solicitudes pueden cambiar tras las revisiones. Los ciclos permiten precisar esos detalles antes de extenderlos a toda la plataforma.
    \item \textbf{La integración puede comprobarse progresivamente.} Se plantea comenzar con consultas de saldos y movimientos para revisar la conexión con el sistema existente, los permisos y la correspondencia de los datos. Después se incorporarían las transferencias internas y las solicitudes de crédito. Cada incremento necesitará pruebas y aprobación antes de ponerse en operación.
    \item \textbf{La disponibilidad del cliente permite revisiones programadas.} La reunión semanal con operaciones puede utilizarse para aclarar dudas y revisar avances. Las decisiones de crédito, tecnología y seguridad se agendarán antes del cierre de cada ciclo. Se supone que estos responsables pueden cumplir esos compromisos; si no ocurre, deberá revisarse la planificación.
    \item \textbf{La documentación puede mantenerse por incremento.} Antes de aceptar cada entrega se actualizarán los requisitos, las interfaces, los criterios de aceptación, la trazabilidad y los resultados de pruebas. Los cambios con impacto en alcance, costo o plazo se someterán a la aprobación prevista en el contrato supuesto.
\end{enumerate}

La matriz de la Tabla~\ref{tab:matriz-cooperativa} respalda esta elección con 17 puntos. La ventaja sobre cascada es reducida y depende principalmente de poder aclarar requisitos e integrar el sistema de cuentas por partes sin abandonar las revisiones formales.

\textbf{Alternativas descartadas para este escenario:}

\textbf{Cascada.} Es una alternativa viable por el control de cambios, la documentación y las aprobaciones que exige la cooperativa. No se elige como organización principal porque todavía quedan requisitos de uso por precisar y conviene comprobar pronto la integración con el sistema existente. Esperar a una etapa posterior para validar estos aspectos aumentaría el retrabajo. Su valoración de 16 puntos muestra que sigue siendo una opción cercana, especialmente si se aclararan los requisitos pendientes antes de construir.

\textbf{Ágil.} También permite revisar requisitos y entregar incrementos, por lo que responde bien a la estabilidad media y al tamaño del proyecto. Se considera menos conveniente en las condiciones asumidas porque socios y especialistas no estarán disponibles de forma continua y deben coordinarse aprobaciones documentales y contractuales. El enfoque ágil puede incorporar esos controles; su descarte no se basa en que carezca de documentación, sino en que la propuesta iterativa planificada se ajusta mejor a la participación programada y a los entregables formales del caso \cite{schon2017,turk2014}.

\subsubsection{Riesgos del modelo elegido y mitigación}

Los siguientes riesgos se plantean a partir de los supuestos de la cooperativa y de la propuesta de desarrollar la plataforma mediante ciclos e incrementos. No representan problemas observados en una organización real.

\begin{enumerate}
    \item \textbf{Retraso de las revisiones.} Si operaciones, crédito o seguridad no revisan un incremento a tiempo, el siguiente ciclo podría comenzar con requisitos pendientes y generar retrabajo. \textbf{Mitigación:} acordar las fechas de revisión al inicio de cada ciclo, enviar previamente los puntos por resolver y registrar las decisiones. No iniciar funciones que dependan de una aprobación pendiente; continuar con tareas que ya tengan requisitos acordados.

    \item \textbf{Crecimiento del alcance entre ciclos.} Las revisiones de pantallas y solicitudes pueden generar nuevas funciones que superen el presupuesto y la fecha pactados. \textbf{Mitigación:} distinguir las aclaraciones de los requisitos de las nuevas solicitudes, estimar su impacto y obtener aprobación antes de incorporarlas. Mantener fuera del alcance inicial las transferencias externas, la aplicación móvil y la aprobación automática de créditos, salvo cambio contractual autorizado.

    \item \textbf{Errores de integración y regresión.} Un incremento de transferencias podría alterar las consultas ya validadas o registrar una operación de forma duplicada al comunicarse con el sistema de cuentas. \textbf{Mitigación:} comprobar la interfaz desde el primer incremento en un entorno de pruebas y repetir las pruebas de las funciones anteriores en cada ciclo. Incluir casos de interrupción, reintento y respuesta fallida; verificar que un reintento no duplique la operación y que saldos y movimientos coincidan con el sistema existente.

    \item \textbf{Controles de seguridad incompletos en entregas parciales.} Al priorizar funciones visibles, el equipo podría posponer permisos o registros necesarios para proteger las consultas y transferencias. \textbf{Mitigación:} incluir requisitos y criterios de aceptación de acceso y registro de operaciones desde el primer ciclo. El responsable de seguridad deberá revisar las evidencias antes de autorizar cada entrega. Una versión mostrada para recibir comentarios no se considerará automáticamente aprobada para operar con datos o dinero reales.

    \item \textbf{Documentación desactualizada.} Los ajustes de cada iteración pueden dejar diferencias entre lo implementado, los requisitos aprobados y las pruebas, dificultando la aceptación del entregable. \textbf{Mitigación:} actualizar la especificación, las interfaces, la trazabilidad y los resultados de pruebas dentro del mismo ciclo. Comprobar su correspondencia con la versión entregada y registrar la aprobación antes de cerrar el incremento.
\end{enumerate}


\section{Matriz comparativa escenario por modelo}
% Ronda 2 (sábado, hasta 18:00)
% Danela crea la tabla. Después, cada integrante llena SOLO la fila de su escenario:
% GAD (Danela), sistema embebido (Tahiny), app móvil (Anthony), cooperativa (José).


\section{Framework de Pohl aplicado al escenario del GAD}
% Responsable: Danela · Ronda 3 (sábado, hasta 22:30)
% Actividades centrales (obtención, documentación, negociación) y transversales
% (validación, gestión) del framework de Pohl en cada fase del modelo elegido para el GAD.


\printbibliography[title={Referencias}]
\end{document}
